\documentclass[11pt]{article}
\usepackage[T1]{fontenc}
\usepackage[utf8]{inputenc}
\usepackage{lmodern}
\usepackage{microtype}
\usepackage[a4paper,margin=27mm]{geometry}
\usepackage{amsmath,amssymb,amsthm,mathtools}
\usepackage{aliascnt}
\usepackage{booktabs,longtable,array,tabularx}
\usepackage{enumitem}
\usepackage{xcolor}
\usepackage{listings}
\usepackage{xurl}
\usepackage{seqsplit}
\usepackage{hyperref}
\usepackage[nameinlink,noabbrev]{cleveref}
\usepackage[backend=biber,style=numeric,sorting=none,maxbibnames=20]{biblatex}
\addbibresource{references.bib}

\hypersetup{
  colorlinks=true,
  linkcolor=blue!45!black,
  citecolor=blue!45!black,
  urlcolor=blue!55!black,
  pdftitle={A 29-AND Circuit for the AES S-box},
  pdfauthor={{GPT-5.6 Pro, umizame} (unordered)},
  pdfsubject={A constructive upper bound on the multiplicative complexity of the AES S-box}
}

\definecolor{codebg}{RGB}{247,247,247}
\definecolor{codeframe}{RGB}{205,205,205}
\lstset{
  basicstyle=\ttfamily\scriptsize,
  backgroundcolor=\color{codebg},
  frame=single,
  rulecolor=\color{codeframe},
  breaklines=true,
  columns=fullflexible,
  keepspaces=true,
  showstringspaces=false,
  aboveskip=0.8em,
  belowskip=0.8em
}

\newtheorem{theorem}{Theorem}
\newaliascnt{proposition}{theorem}
\newtheorem{proposition}[proposition]{Proposition}
\aliascntresetthe{proposition}
\newaliascnt{lemma}{theorem}
\newtheorem{lemma}[lemma]{Lemma}
\aliascntresetthe{lemma}
\newaliascnt{corollary}{theorem}
\newtheorem{corollary}[corollary]{Corollary}
\aliascntresetthe{corollary}
\theoremstyle{definition}
\newaliascnt{definition}{theorem}
\newtheorem{definition}[definition]{Definition}
\aliascntresetthe{definition}
\newaliascnt{remark}{theorem}
\newtheorem{remark}[remark]{Remark}
\aliascntresetthe{remark}

\newcommand{\F}{\mathbb F}
\newcommand{\AES}{\operatorname{AES\mbox{-}Sbox}}
\newcommand{\MC}{\operatorname{MC}}
\newcommand{\Span}{\operatorname{span}}
\newcommand{\sig}{\operatorname{sig}}
\newcommand{\mask}[2]{\left\langle #1\right\rangle_{#2}}
\newcommand{\primaryhash}{\nolinkurl{f41861c4b15bc78840708a5da2fab6ca9dd204c1e3e572cae49912c713bf18d3}}
\newcommand{\transparenthash}{\nolinkurl{cecd0c7abba9920ff2af7ec204fcd01f97db135a4c43d8afd7f83d65c2e382ca}}
\newcommand{\xaghash}{\nolinkurl{862446d12293b21bfafbd7e958502301619a596e6bf3cb7a3b22f3f830926264}}
\newcommand{\towerhash}{\nolinkurl{1e847a85e3abe783d98e3d4c48a7d9eb845caab96903c175657cce06119c6925}}
\newcommand{\stagedhash}{\nolinkurl{d72ca4cd9617b365587a3b826ab3a0d37489187cdd67fca18ff50b1d3d5e0083}}
\newcommand{\baselinehash}{\nolinkurl{6ba4f63b832a1f4520a76cab5680e673fdd744ad2651918882eeecff7957698a}}

\title{\bfseries A 29-AND Circuit for the AES S-box}
\author{\textbf{Authors (unordered):} \{GPT-5.6 Pro, \href{https://github.com/umizame}{umizame}\}}
\date{12 July 2026}

\begin{document}
\maketitle

\begin{abstract}
The NIST Circuit Complexity project defines multiplicative complexity in a one-bit Boolean-circuit
model and lists a forward AES S-box circuit with 32 AND gates.  We give a completely explicit
construction with 29 AND gates.  The construction is specified by eight input-linear coordinates,
24 products of linear forms, five intervening products whose factors are given recursively, four
post-middle coordinates, and eight affine output expressions.  Nine of the 24 outer products are
available before the five middle products and fifteen are available afterward, giving the legal
nonlinear schedule
\[
  9+5+15=29.
\]
The complete formulas and the complete 256-entry evaluation table are included.  The latter agrees
entry by entry with the forward S-box in FIPS 197.  Hence the construction proves
\(
  \MC(\AES)\leq 29.
\)
A concrete straight-line implementation has 29 AND, 195 XOR, and 4 NOT instructions, ordinary gate
depth 35, and AND-depth 6.  Independent verifiers check the staged formulas, the FIPS algebraic
definition and table, the implementation grammar and dependencies, a normalized XOR--AND graph,
the derivation from the official 32-AND NIST program, and deterministic reconstruction.  The result
is an upper bound only; no minimality claim is made.
\end{abstract}

\tableofcontents
\bigskip

\section{Problem statement and exact conventions}
\label{sec:problem}

\subsection{NIST circuit model and multiplicative complexity}

NIST defines the multiplicative size of a Boolean circuit as the number of its nonlinear gates and
the multiplicative complexity of a Boolean or vectorial Boolean function as the minimum
multiplicative size among circuits computing it.  The NIST circuit collection uses the one-bit basis
\(\{\mathrm{AND},\mathrm{XOR},\mathrm{NOT}\}\) for this problem
\cite{nistCircuitComplexity}.  Over \(\F_2\), XOR and NOT are affine and AND is nonlinear; therefore
multiplicative size in this basis is exactly the number of AND gates.

For the forward AES S-box, NIST lists a circuit with 32 AND gates and 81 XOR/XNOR gates, 113 gates in
total, ordinary depth 27, and AND-depth 6 \cite{nistCircuitList}.  The corresponding official
straight-line program contains 32 AND, 77 XOR, and 4 XNOR instructions
\cite{nistAES32SLP}.  NIST asks whether fewer than 32 nonlinear gates suffice
\cite{nistCircuitComplexity}.  The theorem proved below supplies a 29-AND upper bound.

\begin{definition}[Circuit model]
\label{def:circuit-model}
A circuit in this paper has eight input bits, arbitrary fan-out, and one-bit gates of the forms
\[
 a\oplus b,\qquad a\wedge b,\qquad 1\oplus a.
\]
Every gate input must be either an external input or the output of an earlier gate.  The
multiplicative size is the number of \(\wedge\)-gates.  The multiplicative complexity
\(\MC(F)\) is the minimum multiplicative size over all such circuits computing \(F\).
\end{definition}

\subsection{The FIPS 197 function}

Write an external byte as
\[
  x=(U_0U_1\cdots U_7)_2=\sum_{j=0}^{7}U_j2^{7-j},
\]
so \(U_0\) is the most significant bit.  For the polynomial representation, write instead
\(x=\sum_{i=0}^{7}b_i2^i\) and associate the polynomial
\(\sum_{i=0}^{7}b_iX^i\).  FIPS 197 defines the forward S-box by inversion in
\begin{equation}
 \F_2[X]/(X^8+X^4+X^3+X+1),
 \label{eq:aes-field}
\end{equation}
with zero mapped to zero, followed by an affine map with constant byte \texttt{63}
\cite{fips197}.  If \(y_i\) are the least-significant-bit-first coefficients of the inverse, then
\begin{equation}
 s_i=y_i\oplus y_{i+4}\oplus y_{i+5}\oplus y_{i+6}\oplus y_{i+7}\oplus c_i,
 \qquad c=\texttt{63},
 \label{eq:fips-affine}
\end{equation}
where subscripts are reduced modulo 8.  The external output notation is
\(\AES(x)=(S_0S_1\cdots S_7)_2\), with \(S_0=s_7\) and \(S_7=s_0\).
FIPS 197 also prints all 256 output bytes in its S-box table.  Both descriptions are checked
independently in the verification described below.

\subsection{Hexadecimal mask notation}

\begin{definition}[Affine masks]
\label{def:masks}
For an ordered tuple of bits \(B=(b_0,\ldots,b_{r-1})\) and an integer
\(m=\sum_{j=0}^{r-1}m_j2^j\), define
\begin{equation}
  \mask{m}{B}=\bigoplus_{j:m_j=1}b_j.
  \label{eq:mask}
\end{equation}
Masks are printed in hexadecimal.  Thus the least significant bit of the \emph{mask} selects the
first member of the declared tuple; this convention is independent of whether a byte bit is most or
least significant.
\end{definition}

For example, with \(B_U=(U_0,\ldots,U_7)\),
\[
 \mask{\mathtt{69}}{B_U}=U_0\oplus U_3\oplus U_5\oplus U_6.
\]
This convention makes every formula below unambiguous and directly machine-checkable.

\section{The explicit staged 29-AND construction}
\label{sec:construction}

This section defines a Boolean map \(C_{29}:\F_2^8\to\F_2^8\) without reference to any previous
circuit.  Every affine expression is specified by a mask, and every nonlinear operation is displayed
as one product of two affine expressions.

\subsection{Input-linear coordinates and the 24 outer products}

Let
\[
 B_U=(U_0,U_1,\ldots,U_7).
\]
Define \(L_0,\ldots,L_3,R_0,\ldots,R_3\) by the eight masks in
\cref{tab:coordinates}.  These are linear functions of the input and require no AND gate.

Eventually there will also be four bits \(Z_0,\ldots,Z_3\).  Put
\begin{equation}
 B_C=(L_0,L_1,L_2,L_3,R_0,R_1,R_2,R_3,Z_0,Z_1,Z_2,Z_3).
 \label{eq:coordinate-basis}
\end{equation}
For each row \((u_j,v_j)\) of \cref{tab:products}, define
\begin{equation}
  P_j=\mask{u_j}{B_C}\wedge\mask{v_j}{B_C}.
  \label{eq:outer-products}
\end{equation}
The nine products
\begin{equation}
 E=(P_{01},P_{02},P_{03},P_{09},P_{10},P_{11},P_{17},P_{18},P_{19})
 \label{eq:early-products}
\end{equation}
have zero coefficients on all four \(Z\)-coordinates, as is visible directly from their factor
masks.  They can therefore be computed immediately after \(L\) and \(R\).  The other fifteen
products are deferred until \(Z\) has been defined.

\subsection{Five recursively available middle products}

Use the padded ordered tuple
\begin{equation}
 B_M=(1,U_0,\ldots,U_7,E,M_1,M_2,M_3,M_4,M_5).
 \label{eq:middle-basis}
\end{equation}
For \(k=1,\ldots,5\), the row for \(M_k\) in \cref{tab:middle-factors} contains two masks
\(a_k,b_k\).  Define recursively
\begin{equation}
  M_k=\mask{a_k}{B_M}\wedge\mask{b_k}{B_M}.
  \label{eq:middle-products}
\end{equation}
Every coefficient on \(M_k,M_{k+1},\ldots,M_5\) is zero in that row.  Consequently the right side of
\cref{eq:middle-products} uses only the constant, inputs, the nine already-computed products in
\(E\), and earlier middle values \(M_1,\ldots,M_{k-1}\).  The recursion is therefore a legal
straight-line schedule rather than a simultaneous or circular definition.

\subsection{Post-middle coordinates, late products, and outputs}

After \(M_5\) has been computed, define \(Z_0,\ldots,Z_3\) by the four masks in
\cref{tab:post-coordinates}, all interpreted in \(B_M\).  The full coordinate tuple \(B_C\) in
\cref{eq:coordinate-basis} is then available, so the remaining fifteen rows of
\cref{tab:products} define the late products
\begin{equation}
 D=(P_{04},P_{05},P_{06},P_{07},P_{08},P_{12},P_{13},P_{14},P_{15},P_{16},
 P_{20},P_{21},P_{22},P_{23},P_{24}).
 \label{eq:late-products}
\end{equation}
Finally set
\[
 B_O=(B_M,D)
\]
and define \(S_0,\ldots,S_7\) by the eight masks in
\cref{tab:output-reconstruction}, interpreted in \(B_O\).

\begin{lemma}[Availability and nonlinear count]
\label{lem:availability}
The construction above is an acyclic circuit in the model of \cref{def:circuit-model}.  It contains
exactly 29 AND gates, with schedule \(9+5+15\).
\end{lemma}

\begin{proof}
The eight \(L,R\) coordinates are linear in the inputs.  By inspection of the nine masks listed in
\cref{eq:early-products}, their factors involve only \(L,R\), so those nine products are available
first.  In the row defining \(M_k\), all mask bits belonging to \(M_k\) and later middle values are
zero; induction on \(k\) therefore proves that each of the five middle products is available when it
is formed.  The four \(Z\)-coordinates are affine in values then available.  This makes every factor
of the fifteen late products available.  The outputs are affine combinations of already-computed
values.  No affine step contributes an AND gate.  The only nonlinear operations are the nine early
products, five middle products, and fifteen late products, totaling \(9+5+15=29\).
\end{proof}

\section{Correctness and the multiplicative-complexity bound}
\label{sec:correctness}

The construction has only eight inputs, so functional correctness can be certified without sampling
or an unproved algebraic identification: evaluate all \(2^8=256\) inputs.  \cref{app:finite} gives the complete output table obtained from the recursion in
\cref{eq:outer-products,eq:middle-products}.  Row \(r\) and column \(c\) contain the output for input
byte \(\mathtt{rc}\).  It is entrywise identical to FIPS 197, Table 4.

\begin{lemma}[Finite-domain equality]
\label{lem:finite-domain}
Let \(F,G:\F_2^8\to\F_2^8\).  If \(F(x)=G(x)\) for every one of the 256 input bytes, then
\(F=G\).
\end{lemma}

\begin{proof}
The 256 bytes are exactly all elements of \(\F_2^8\).  Pointwise equality on this complete domain is
the definition of equality of the two functions.
\end{proof}

\begin{theorem}[Constructive 29-AND upper bound]
\label{thm:main}
For the forward AES S-box specified by FIPS 197, in the circuit model of
\cref{def:circuit-model},
\[
  \MC(\AES)\leq 29.
\]
\end{theorem}

\begin{proof}
By \cref{lem:availability}, \(C_{29}\) is a legal circuit with exactly 29 AND gates.  Direct recursive
evaluation of its displayed formulas gives \cref{tab:evaluation}.  The same table is the complete
FIPS 197 forward S-box table.  Therefore \(C_{29}(x)=\AES(x)\) for every input byte; by
\cref{lem:finite-domain}, \(C_{29}=\AES\).  Since multiplicative complexity is the minimum AND count
over all circuits computing the function, the existence of this circuit gives
\(\MC(\AES)\leq29\).
\end{proof}

The table is not the only correctness check.  The standalone verifier reproduced in
\cref{app:artifacts} first checks the literal 256-byte FIPS table against the field inversion and
affine map in \cref{eq:aes-field,eq:fips-affine}; it then parses and exhaustively evaluates the
released straight-line implementation.  A separate staged verifier interprets precisely the mask
bases and recursions in this paper, enforces every availability restriction, and again compares all
256 values with FIPS.  Thus the field definition, printed table, staged construction, and compiled
straight-line program are pairwise connected by explicit finite checks.

\begin{corollary}[Concrete straight-line implementation]
\label{cor:slp}
There is a straight-line implementation of \(C_{29}\) containing
\[
 29\ \mathrm{AND}+195\ \mathrm{XOR}+4\ \mathrm{NOT}=228
\]
instructions.  Its ordinary gate depth is 35 and its AND-depth is 6.
\end{corollary}

\begin{proof}
The strict parser listed in \cref{app:artifacts} accepts only the declared one-bit grammar, verifies
that every operand is previously defined, prohibits wire redefinition, requires all declared
temporaries and outputs, and counts the accepted instructions.  It obtains the stated tally.  Giving
inputs depth zero and assigning each gate one plus the maximum operand depth yields output depth 35.
For AND-depth, the same recurrence adds one only at an AND gate and obtains 6.  Because the verified
program order is acyclic, these recurrences equal the corresponding longest-path definitions.  The
same parser exhaustively verifies the implementation output against FIPS on all 256 inputs.
\end{proof}

\section{Independent derivation from the official NIST circuit}
\label{sec:baseline}

The proof of \cref{thm:main} is already complete and does not assume the correctness of another
circuit.  This section records how the construction was obtained and supplies a second, structurally
different audit.

The official NIST 32-AND SLP numbers its nonlinear outputs in program order.  The first nine are
\[
 t24,t25,t27,t29,t30,t32,t34,t35,t37,
\]
the next five are
\[
 t48,t50,t54,t58,t61,
\]
and the last eighteen are \(t68,t69,\ldots,t85\).  Hence its nonlinear schedule is
\(9+5+18\).  The SLP states that it is based on the low-multiplicative-complexity AES construction of
Boyar, Matthews, and Peralta, with an affine-layer improvement
\cite{nistAES32SLP,boyarMatthewsPeralta2013}.

The following coordinate vectors are exact named-wire identities in that SLP:
\begin{equation}
\begin{aligned}
 L&=(t8,t11,t2,t10),\\
 R&=(t13,t15,t21,t6),\\
 Z&=(t66,t59,t65,t62).
\end{aligned}
\label{eq:LRZ}
\end{equation}
The first two are available before any AND gate and equal the input-linear forms printed in
\cref{tab:coordinates}.  The vector \(Z\) is available after the five middle AND gates.  The new
construction retains those five middle functions, replaces the 27 outer products by the 24 products
of \cref{tab:products}, and supplies explicit affine reconstruction masks for every retained factor,
every \(Z_i\), and every output.  The exact masks, rather than an assertion of span membership, are
printed in \cref{tab:middle-factors,tab:post-coordinates,tab:output-reconstruction}.

For an additional symbolic audit, every NIST wire is processed as an affine expression over the
constant, the eight inputs, and 32 independent symbols representing the successive NIST AND outputs.
XOR, XNOR, and NOT update this expression exactly; encountering an AND introduces the next independent
symbol.  This syntactic calculation does not assume a truth-table identity between unrelated wires.
It proves the coordinate identities in \cref{eq:LRZ}, the factor-coordinate forms of all 27 outer
AND gates, and the exact coefficient vectors by which outer AND outputs enter the ten middle factors
and eight final outputs.

The remaining part of that symbolic bridge is described in \cref{sec:formal}.  Its role is
provenance and redundancy: the direct staged formulas and complete FIPS evaluation remain sufficient
for the theorem even if the NIST derivation is ignored.

\section{An eight-product identity over \texorpdfstring{\(\F_4\)}{F4}}
\label{sec:gf4}

The 24 outer products have a compact algebraic organization.  Let
\[
 K=\F_4=\F_2[\omega]/(\omega^2+\omega+1)
\]
and write \(x=x_0\omega+x_1\), \(y=y_0\omega+y_1\), and
\(z=z_0\omega+z_1\).  Define
\begin{align}
 p_1&=x_0y_0, &
 p_2&=x_1y_1, &
 p_3&=(x_0\oplus x_1)(y_0\oplus y_1),
 \label{eq:p123}\\
 p_4&=(x_0\oplus y_0\oplus y_1\oplus z_1)(x_1\oplus y_0\oplus z_0),\nonumber\\
 p_5&=(x_1\oplus y_0\oplus y_1\oplus z_1)(x_0\oplus x_1\oplus y_0\oplus z_0),\nonumber\\
 p_6&=(x_1\oplus y_0\oplus y_1\oplus z_0\oplus z_1)(x_0\oplus y_1\oplus z_0),\nonumber\\
 p_7&=(x_1\oplus y_0\oplus z_1)(x_0\oplus x_1\oplus y_1\oplus z_0),\nonumber\\
 p_8&=(x_1\oplus y_0\oplus z_0\oplus z_1)(x_0\oplus y_0\oplus y_1\oplus z_0).
 \label{eq:p45678}
\end{align}
Then define six affine reconstructions by
\begin{align}
 (xy)_0&=p_2\oplus p_3, &
 (xy)_1&=p_1\oplus p_2,
 \label{eq:xy}\\
 (yz)_0&=y_0\oplus y_1\oplus p_1\oplus p_2\oplus p_3\oplus p_6\oplus p_8,\nonumber\\
 (yz)_1&=y_1\oplus p_5\oplus p_6\oplus p_7\oplus p_8,
 \label{eq:yz}\\
 (zx)_0&=x_1\oplus p_1\oplus p_2\oplus p_3\oplus p_4\oplus p_5,\nonumber\\
 (zx)_1&=x_1\oplus y_1\oplus z_0\oplus p_1\oplus p_2\oplus p_3\oplus p_6\oplus p_7.
 \label{eq:zx}
\end{align}
Here \((ab)_0\omega+(ab)_1\) denotes the displayed pair of reconstructed bits.

\begin{proposition}[Filtered \(\F_4\) triangle identity]
\label{prop:gf4}
For every \((x,y,z)\in K^3\), the quantities in
\cref{eq:xy,eq:yz,eq:zx} are exactly the coordinates of \(xy,yz,zx\).  Products
\(p_1,p_2,p_3\) use only \(x,y\), while \(p_4,\ldots,p_8\) may use \(z\).
\end{proposition}

\begin{proof}
Since \(\omega^2=\omega+1\),
\begin{equation}
 (a_0\omega+a_1)(b_0\omega+b_1)
 =(a_0b_0\oplus a_0b_1\oplus a_1b_0)\omega
  \oplus(a_0b_0\oplus a_1b_1).
 \label{eq:gf4mul}
\end{equation}
Expand the right sides of \cref{eq:xy,eq:yz,eq:zx}, substitute
\cref{eq:p123,eq:p45678}, and cancel repeated monomials over \(\F_2\).  The six resulting algebraic
normal forms are
\begin{align*}
 (xy)_0&=x_0y_0\oplus x_0y_1\oplus x_1y_0,
 & (xy)_1&=x_0y_0\oplus x_1y_1,\\
 (yz)_0&=y_0z_0\oplus y_0z_1\oplus y_1z_0,
 & (yz)_1&=y_0z_0\oplus y_1z_1,\\
 (zx)_0&=z_0x_0\oplus z_0x_1\oplus z_1x_0,
 & (zx)_1&=z_0x_0\oplus z_1x_1.
\end{align*}
These are precisely the two coordinates in \cref{eq:gf4mul}.  The availability statement follows
directly from the eight factor formulas.
\end{proof}

To obtain the 24 rows in \cref{tab:products}, first apply the invertible binary transform
\begin{equation}
 \Phi(a_0,a_1,a_2,a_3)=
 \bigl((a_2\oplus a_3)\omega+a_3,\ a_1\omega+a_0\bigr).
 \label{eq:Phi}
\end{equation}
Its four binary row masks are \(\mathtt c,\mathtt8,\mathtt2,\mathtt1\); the compact tag
\(\mathtt{128c}\) records them from high to low hexadecimal digit.  Invertibility is immediate from
\[
 a_3=q_1,\qquad a_2=q_0\oplus q_1,\qquad a_1=q_2,\qquad a_0=q_3,
\]
where \((q_0,q_1,q_2,q_3)=(a_2\oplus a_3,a_3,a_1,a_0)\).

Apply \cref{prop:gf4} three times: to the first components of
\(\Phi(L),\Phi(R),\Phi(Z)\), to their second components, and to the componentwise sums of those two
triples.  This produces three blocks of eight products.  In each block, three products are early and
five are late, so the total filtration is nine early and fifteen late.  This observation explains the
factor pattern in \cref{tab:products}.  It is not an assumption that a particular field structure on
four-bit vectors matches the NIST circuit: the exact relation to the required NIST functions is
checked by the explicit reconstruction masks and the formal bridge below.

\section{Formal quadratic bridge}
\label{sec:formal}

The following calculation independently checks why the 24 factor pairs can replace the outer portion
of the 32-AND NIST construction.

Let \(V=\F_2^{12}\) with ordered coordinates \(B_C\) from
\cref{eq:coordinate-basis}.  For \(u\in V\), let \(\ell_u=\mask{u}{B_C}\), now treating the twelve
coordinates as independent Boolean indeterminates.

\begin{lemma}[Product signature and diagonal correction]
\label{lem:wedge}
In the Boolean polynomial ring with \(X_i^2=X_i\),
\begin{equation}
 \ell_u(X)\ell_v(X)=
 \bigoplus_{i<j}(u_iv_j\oplus u_jv_i)X_iX_j
 \oplus\bigoplus_i u_iv_iX_i.
 \label{eq:wedge}
\end{equation}
Thus the off-diagonal quadratic signature is the decomposable two-form
\(u\wedge v\in\Lambda^2V\), while the diagonal linear correction has mask
\(u\mathbin{\&}v\).
\end{lemma}

\begin{proof}
Expand the product.  For \(i<j\), the coefficient of \(X_iX_j\) is the XOR of the two ordered
contributions \(u_iv_j\) and \(u_jv_i\).  The diagonal contribution is
\(u_iv_iX_i^2=u_iv_iX_i\).  Grouping the terms gives \cref{eq:wedge}.
\end{proof}

Encode \(u\wedge v\) in the lexicographically ordered pair basis
\[
 (0,1),(0,2),\ldots,(0,11),(1,2),\ldots,(10,11).
\]
A symbolic pass through the NIST SLP gives 27 linearly independent signatures for its nine early and
18 late outer AND outputs.  Projecting the ten middle-factor expressions and eight output expressions
onto those 27 outer symbols, and then mapping through the corresponding signatures, gives an ordered
four-dimensional pre-middle space \(Q_{\mathrm{pre}}\) and an ordered twelve-dimensional total space
\(Q\), with
\[
 Q_{\mathrm{pre}}\subset Q\subset\Lambda^2V,
 \qquad \dim Q_{\mathrm{pre}}=4,
 \qquad \dim Q=12.
\]
The twelve target signatures and the 24-bit product-selection masks are printed in
\cref{tab:targets}; the first four selection masks use only the nine early products.

\begin{proposition}[Exact formal reconstruction]
\label{prop:formal-certificate}
For every row of \cref{tab:targets}, the selected products from \cref{tab:products}, together with
the diagonal correction prescribed by \cref{lem:wedge}, equal the indicated quadratic polynomial on
all \(2^{12}=4096\) assignments of the formal coordinates.  The first four identities use no late
product.
\end{proposition}

\begin{proof}
For a selection set \(J\), compute
\[
 d=\bigoplus_{j\in J}(u_j\wedge v_j),
 \qquad
 r=\bigoplus_{j\in J}(u_j\mathbin{\&}v_j).
\]
The stored selection mask is checked to give exactly the target signature \(d\).  Applying
\cref{lem:wedge} term by term then gives the polynomial identity
\[
 q(X)=\bigoplus_{j\in J}\ell_{u_j}(X)\ell_{v_j}(X)\oplus\ell_r(X).
\]
This coefficient argument proves the identity.  As a redundant check of coordinate and bit ordering,
both sides are also evaluated on all 4096 assignments.  Inspection of the first four selection masks
shows that their set bits occur only at the nine early product positions.
\end{proof}

The formal calculation concerns the quadratic outer contribution.  The full affine corrections,
constant terms, recursive middle dependencies, post-middle coordinates, and output expressions are
not left implicit: they are the explicit staged masks in \cref{tab:middle-factors,tab:post-coordinates,tab:output-reconstruction}.
Their exhaustive eight-input evaluation is the direct correctness proof used in
\cref{thm:main}.

\section{Verification separation}
\label{sec:verification}

The result is checked through independent representations so that a defect in one representation is
unlikely to be silently inherited by all others.

\begin{enumerate}[leftmargin=2.4em]
\item The staged-certificate checker parses the 24 factor pairs and all affine masks, enforces the
nine-before-five-before-fifteen availability filtration, and evaluates all 256 inputs against a
literal FIPS table.
\item The direct straight-line checker strictly parses declarations, delimiters, gate arities,
dependencies, and wire names; it recomputes the exact tally and both depths, independently derives the
FIPS S-box by finite-field arithmetic, checks that derivation against the literal table, and evaluates
the circuit on all inputs.
\item A bit-parallel evaluator represents every wire by its complete 256-bit truth table.  A separate
C99 parser and evaluator repeats the finite circuit test.
\item Both straight-line implementations are normalized by absorbing every affine operation into
masks over the constant, inputs, and preceding AND outputs.  Their 29 ordered factor pairs and eight
output masks agree exactly with the printed XOR--AND graph certificate.
\item The NIST bridge derives its target from the official 32-AND SLP rather than accepting a target
from the 24-product witness; it verifies the formal identities on all 4096 coordinate assignments and
the actual staged realization on all 256 AES inputs.
\item A clean-tree rebuild regenerates the readable circuit, resynthesizes only its affine network,
regenerates the normalized graph, and compares all generated artifacts byte for byte.
\end{enumerate}

These checks have deliberately different proof obligations.  The first two alone each give a
complete finite correctness check for the upper bound.  The remaining checks establish provenance,
representation consistency, and deterministic reproducibility.  None of them is used to assert a
29-gate lower bound.

\section{Conclusion}
\label{sec:conclusion}

The explicit staged construction in this paper is a legal one-bit Boolean circuit with nonlinear
schedule \(9+5+15=29\).  Its complete 256-entry output table is the FIPS 197 forward AES S-box table.
Therefore \(\MC(\AES)\leq29\).  A concrete implementation uses 228 one-bit instructions and has
AND-depth 6.  The derivation from the official NIST 32-AND construction, the three-block
\(\F_4\) identity, the formal quadratic bridge, and multiple independent evaluators provide
additional checks, but the theorem ultimately rests on the explicit formulas and complete finite
evaluation.  The construction establishes an upper bound and does not establish optimality.

\appendix

\section{Complete staged decomposition}
\label{app:decomposition}

All masks in this appendix use \cref{def:masks}.  The product table supplies every factor of all 24
outer AND gates.  The staged tables then supply every input coordinate, every factor of the five
middle AND gates, every post-middle coordinate, and every output expression.  Together these tables
are a complete specification of \(C_{29}\).

\begingroup\small
\begin{longtable}{@{}rcc>{\raggedright\arraybackslash}p{0.31\textwidth}>{\raggedright\arraybackslash}p{0.31\textwidth}@{}}
\caption{Complete 24-product certificate. A factor mask uses coordinate order $(L_0,\ldots,L_3,R_0,\ldots,R_3,Z_0,\ldots,Z_3)$. Multiplication is commutative, so the displayed factor order has no mathematical significance.}\label{tab:products}\\
\toprule
product & stage & block & first factor & second factor \\
\midrule
\endfirsthead
\toprule
product & stage & block & first factor & second factor \\
\midrule
\endhead
\bottomrule
\endfoot
$P_{01}$ & early & $0$ & $\ell_{\mathtt{0c0}}=R_2\mathbin{\oplus}\allowbreak R_3$ & $\ell_{\mathtt{00c}}=L_2\mathbin{\oplus}\allowbreak L_3$ \\
$P_{02}$ & early & $0$ & $\ell_{\mathtt{080}}=R_3$ & $\ell_{\mathtt{008}}=L_3$ \\
$P_{03}$ & early & $0$ & $\ell_{\mathtt{040}}=R_2$ & $\ell_{\mathtt{004}}=L_2$ \\
$P_{04}$ & late & $0$ & $\ell_{\mathtt{84c}}=L_2\mathbin{\oplus}\allowbreak L_3\mathbin{\oplus}\allowbreak R_2\mathbin{\oplus}\allowbreak Z_3$ & $\ell_{\mathtt{cc8}}=L_3\mathbin{\oplus}\allowbreak R_2\mathbin{\oplus}\allowbreak R_3\mathbin{\oplus}\allowbreak Z_2\mathbin{\oplus}\allowbreak Z_3$ \\
$P_{05}$ & late & $0$ & $\ell_{\mathtt{848}}=L_3\mathbin{\oplus}\allowbreak R_2\mathbin{\oplus}\allowbreak Z_3$ & $\ell_{\mathtt{cc4}}=L_2\mathbin{\oplus}\allowbreak R_2\mathbin{\oplus}\allowbreak R_3\mathbin{\oplus}\allowbreak Z_2\mathbin{\oplus}\allowbreak Z_3$ \\
$P_{06}$ & late & $0$ & $\ell_{\mathtt{448}}=L_3\mathbin{\oplus}\allowbreak R_2\mathbin{\oplus}\allowbreak Z_2$ & $\ell_{\mathtt{c8c}}=L_2\mathbin{\oplus}\allowbreak L_3\mathbin{\oplus}\allowbreak R_3\mathbin{\oplus}\allowbreak Z_2\mathbin{\oplus}\allowbreak Z_3$ \\
$P_{07}$ & late & $0$ & $\ell_{\mathtt{8c8}}=L_3\mathbin{\oplus}\allowbreak R_2\mathbin{\oplus}\allowbreak R_3\mathbin{\oplus}\allowbreak Z_3$ & $\ell_{\mathtt{c84}}=L_2\mathbin{\oplus}\allowbreak R_3\mathbin{\oplus}\allowbreak Z_2\mathbin{\oplus}\allowbreak Z_3$ \\
$P_{08}$ & late & $0$ & $\ell_{\mathtt{4c8}}=L_3\mathbin{\oplus}\allowbreak R_2\mathbin{\oplus}\allowbreak R_3\mathbin{\oplus}\allowbreak Z_2$ & $\ell_{\mathtt{c4c}}=L_2\mathbin{\oplus}\allowbreak L_3\mathbin{\oplus}\allowbreak R_2\mathbin{\oplus}\allowbreak Z_2\mathbin{\oplus}\allowbreak Z_3$ \\
$P_{09}$ & early & $1$ & $\ell_{\mathtt{020}}=R_1$ & $\ell_{\mathtt{002}}=L_1$ \\
$P_{10}$ & early & $1$ & $\ell_{\mathtt{010}}=R_0$ & $\ell_{\mathtt{001}}=L_0$ \\
$P_{11}$ & early & $1$ & $\ell_{\mathtt{030}}=R_0\mathbin{\oplus}\allowbreak R_1$ & $\ell_{\mathtt{003}}=L_0\mathbin{\oplus}\allowbreak L_1$ \\
$P_{12}$ & late & $1$ & $\ell_{\mathtt{132}}=L_1\mathbin{\oplus}\allowbreak R_0\mathbin{\oplus}\allowbreak R_1\mathbin{\oplus}\allowbreak Z_0$ & $\ell_{\mathtt{221}}=L_0\mathbin{\oplus}\allowbreak R_1\mathbin{\oplus}\allowbreak Z_1$ \\
$P_{13}$ & late & $1$ & $\ell_{\mathtt{131}}=L_0\mathbin{\oplus}\allowbreak R_0\mathbin{\oplus}\allowbreak R_1\mathbin{\oplus}\allowbreak Z_0$ & $\ell_{\mathtt{223}}=L_0\mathbin{\oplus}\allowbreak L_1\mathbin{\oplus}\allowbreak R_1\mathbin{\oplus}\allowbreak Z_1$ \\
$P_{14}$ & late & $1$ & $\ell_{\mathtt{331}}=L_0\mathbin{\oplus}\allowbreak R_0\mathbin{\oplus}\allowbreak R_1\mathbin{\oplus}\allowbreak Z_0\mathbin{\oplus}\allowbreak Z_1$ & $\ell_{\mathtt{212}}=L_1\mathbin{\oplus}\allowbreak R_0\mathbin{\oplus}\allowbreak Z_1$ \\
$P_{15}$ & late & $1$ & $\ell_{\mathtt{121}}=L_0\mathbin{\oplus}\allowbreak R_1\mathbin{\oplus}\allowbreak Z_0$ & $\ell_{\mathtt{213}}=L_0\mathbin{\oplus}\allowbreak L_1\mathbin{\oplus}\allowbreak R_0\mathbin{\oplus}\allowbreak Z_1$ \\
$P_{16}$ & late & $1$ & $\ell_{\mathtt{321}}=L_0\mathbin{\oplus}\allowbreak R_1\mathbin{\oplus}\allowbreak Z_0\mathbin{\oplus}\allowbreak Z_1$ & $\ell_{\mathtt{232}}=L_1\mathbin{\oplus}\allowbreak R_0\mathbin{\oplus}\allowbreak R_1\mathbin{\oplus}\allowbreak Z_1$ \\
$P_{17}$ & early & $0\mathbin{\oplus}1$ & $\ell_{\mathtt{0e0}}=R_1\mathbin{\oplus}\allowbreak R_2\mathbin{\oplus}\allowbreak R_3$ & $\ell_{\mathtt{00e}}=L_1\mathbin{\oplus}\allowbreak L_2\mathbin{\oplus}\allowbreak L_3$ \\
$P_{18}$ & early & $0\mathbin{\oplus}1$ & $\ell_{\mathtt{090}}=R_0\mathbin{\oplus}\allowbreak R_3$ & $\ell_{\mathtt{009}}=L_0\mathbin{\oplus}\allowbreak L_3$ \\
$P_{19}$ & early & $0\mathbin{\oplus}1$ & $\ell_{\mathtt{070}}=R_0\mathbin{\oplus}\allowbreak R_1\mathbin{\oplus}\allowbreak R_2$ & $\ell_{\mathtt{007}}=L_0\mathbin{\oplus}\allowbreak L_1\mathbin{\oplus}\allowbreak L_2$ \\
$P_{20}$ & late & $0\mathbin{\oplus}1$ & $\ell_{\mathtt{97e}}=L_1\mathbin{\oplus}\allowbreak L_2\mathbin{\oplus}\allowbreak L_3\mathbin{\oplus}\allowbreak R_0\mathbin{\oplus}\allowbreak R_1\mathbin{\oplus}\allowbreak R_2\mathbin{\oplus}\allowbreak Z_0\mathbin{\oplus}\allowbreak Z_3$ & $\ell_{\mathtt{ee9}}=L_0\mathbin{\oplus}\allowbreak L_3\mathbin{\oplus}\allowbreak R_1\mathbin{\oplus}\allowbreak R_2\mathbin{\oplus}\allowbreak R_3\mathbin{\oplus}\allowbreak Z_1\mathbin{\oplus}\allowbreak Z_2\mathbin{\oplus}\allowbreak Z_3$ \\
$P_{21}$ & late & $0\mathbin{\oplus}1$ & $\ell_{\mathtt{979}}=L_0\mathbin{\oplus}\allowbreak L_3\mathbin{\oplus}\allowbreak R_0\mathbin{\oplus}\allowbreak R_1\mathbin{\oplus}\allowbreak R_2\mathbin{\oplus}\allowbreak Z_0\mathbin{\oplus}\allowbreak Z_3$ & $\ell_{\mathtt{ee7}}=L_0\mathbin{\oplus}\allowbreak L_1\mathbin{\oplus}\allowbreak L_2\mathbin{\oplus}\allowbreak R_1\mathbin{\oplus}\allowbreak R_2\mathbin{\oplus}\allowbreak R_3\mathbin{\oplus}\allowbreak Z_1\mathbin{\oplus}\allowbreak Z_2\mathbin{\oplus}\allowbreak Z_3$ \\
$P_{22}$ & late & $0\mathbin{\oplus}1$ & $\ell_{\mathtt{779}}=L_0\mathbin{\oplus}\allowbreak L_3\mathbin{\oplus}\allowbreak R_0\mathbin{\oplus}\allowbreak R_1\mathbin{\oplus}\allowbreak R_2\mathbin{\oplus}\allowbreak Z_0\mathbin{\oplus}\allowbreak Z_1\mathbin{\oplus}\allowbreak Z_2$ & $\ell_{\mathtt{e9e}}=L_1\mathbin{\oplus}\allowbreak L_2\mathbin{\oplus}\allowbreak L_3\mathbin{\oplus}\allowbreak R_0\mathbin{\oplus}\allowbreak R_3\mathbin{\oplus}\allowbreak Z_1\mathbin{\oplus}\allowbreak Z_2\mathbin{\oplus}\allowbreak Z_3$ \\
$P_{23}$ & late & $0\mathbin{\oplus}1$ & $\ell_{\mathtt{9e9}}=L_0\mathbin{\oplus}\allowbreak L_3\mathbin{\oplus}\allowbreak R_1\mathbin{\oplus}\allowbreak R_2\mathbin{\oplus}\allowbreak R_3\mathbin{\oplus}\allowbreak Z_0\mathbin{\oplus}\allowbreak Z_3$ & $\ell_{\mathtt{e97}}=L_0\mathbin{\oplus}\allowbreak L_1\mathbin{\oplus}\allowbreak L_2\mathbin{\oplus}\allowbreak R_0\mathbin{\oplus}\allowbreak R_3\mathbin{\oplus}\allowbreak Z_1\mathbin{\oplus}\allowbreak Z_2\mathbin{\oplus}\allowbreak Z_3$ \\
$P_{24}$ & late & $0\mathbin{\oplus}1$ & $\ell_{\mathtt{7e9}}=L_0\mathbin{\oplus}\allowbreak L_3\mathbin{\oplus}\allowbreak R_1\mathbin{\oplus}\allowbreak R_2\mathbin{\oplus}\allowbreak R_3\mathbin{\oplus}\allowbreak Z_0\mathbin{\oplus}\allowbreak Z_1\mathbin{\oplus}\allowbreak Z_2$ & $\ell_{\mathtt{e7e}}=L_1\mathbin{\oplus}\allowbreak L_2\mathbin{\oplus}\allowbreak L_3\mathbin{\oplus}\allowbreak R_0\mathbin{\oplus}\allowbreak R_1\mathbin{\oplus}\allowbreak R_2\mathbin{\oplus}\allowbreak Z_1\mathbin{\oplus}\allowbreak Z_2\mathbin{\oplus}\allowbreak Z_3$ \\
\end{longtable}
\endgroup

\subsection{Input-linear coordinates}
The input bits are ordered $(U_0,\ldots,U_7)$, with $U_0$ the most significant bit.
\begin{longtable}{@{}rcl@{}}
\caption{The eight pre-middle coordinates as linear forms in $(U_0,\ldots,U_7)$.}\label{tab:coordinates}\\
\toprule
coordinate & mask & linear form \\
\midrule\endfirsthead
\toprule
coordinate & mask & linear form \\
\midrule\endhead
\bottomrule\endfoot
$L_0$ & \texttt{69} & $U_0\mathbin{\oplus}\allowbreak U_3\mathbin{\oplus}\allowbreak U_5\mathbin{\oplus}\allowbreak U_6$ \\
$L_1$ & \texttt{e7} & $U_0\mathbin{\oplus}\allowbreak U_1\mathbin{\oplus}\allowbreak U_2\mathbin{\oplus}\allowbreak U_5\mathbin{\oplus}\allowbreak U_6\mathbin{\oplus}\allowbreak U_7$ \\
$L_2$ & \texttt{41} & $U_0\mathbin{\oplus}\allowbreak U_6$ \\
$L_3$ & \texttt{c6} & $U_1\mathbin{\oplus}\allowbreak U_2\mathbin{\oplus}\allowbreak U_6\mathbin{\oplus}\allowbreak U_7$ \\
$R_0$ & \texttt{59} & $U_0\mathbin{\oplus}\allowbreak U_3\mathbin{\oplus}\allowbreak U_4\mathbin{\oplus}\allowbreak U_6$ \\
$R_1$ & \texttt{d9} & $U_0\mathbin{\oplus}\allowbreak U_3\mathbin{\oplus}\allowbreak U_4\mathbin{\oplus}\allowbreak U_6\mathbin{\oplus}\allowbreak U_7$ \\
$R_2$ & \texttt{74} & $U_2\mathbin{\oplus}\allowbreak U_4\mathbin{\oplus}\allowbreak U_5\mathbin{\oplus}\allowbreak U_6$ \\
$R_3$ & \texttt{86} & $U_1\mathbin{\oplus}\allowbreak U_2\mathbin{\oplus}\allowbreak U_7$ \\
\end{longtable}

\subsection{Five retained middle products}
For this table the mask basis is
\[
(1,U_0,\ldots,U_7,P_{01},P_{02},P_{03},P_{09},P_{10},P_{11},P_{17},P_{18},P_{19},M_1,\ldots,M_5).
\]
A row defining $M_k$ has zero coefficients on $M_k,\ldots,M_5$; hence both factors are available before $M_k$ is formed.
\begingroup\small
\begin{longtable}{@{}r>{\raggedright\arraybackslash}p{0.41\textwidth}>{\raggedright\arraybackslash}p{0.41\textwidth}@{}}
\caption{Exact affine factors of the five middle AND gates.}\label{tab:middle-factors}\\
\toprule
gate & first factor & second factor \\
\midrule\endfirsthead
\toprule
gate & first factor & second factor \\
\midrule\endhead
\bottomrule\endfoot
$M_1$ & $\ell_{\mathtt{1ecf6}}=U_0\mathbin{\oplus}\allowbreak U_1\mathbin{\oplus}\allowbreak U_3\mathbin{\oplus}\allowbreak U_4\mathbin{\oplus}\allowbreak U_5\mathbin{\oplus}\allowbreak U_6\mathbin{\oplus}\allowbreak P_{02}\mathbin{\oplus}\allowbreak P_{03}\mathbin{\oplus}\allowbreak P_{10}\mathbin{\oplus}\allowbreak P_{11}\mathbin{\oplus}\allowbreak P_{17}\mathbin{\oplus}\allowbreak P_{18}$ & $\ell_{\mathtt{1d06a}}=U_0\mathbin{\oplus}\allowbreak U_2\mathbin{\oplus}\allowbreak U_4\mathbin{\oplus}\allowbreak U_5\mathbin{\oplus}\allowbreak P_{09}\mathbin{\oplus}\allowbreak P_{11}\mathbin{\oplus}\allowbreak P_{17}\mathbin{\oplus}\allowbreak P_{18}$ \\
$M_2$ & $\ell_{\mathtt{2b60a}}=U_0\mathbin{\oplus}\allowbreak U_2\mathbin{\oplus}\allowbreak P_{01}\mathbin{\oplus}\allowbreak P_{02}\mathbin{\oplus}\allowbreak P_{09}\mathbin{\oplus}\allowbreak P_{10}\mathbin{\oplus}\allowbreak P_{17}\mathbin{\oplus}\allowbreak P_{19}$ & $\ell_{\mathtt{730ea}}=U_0\mathbin{\oplus}\allowbreak U_2\mathbin{\oplus}\allowbreak U_4\mathbin{\oplus}\allowbreak U_5\mathbin{\oplus}\allowbreak U_6\mathbin{\oplus}\allowbreak P_{09}\mathbin{\oplus}\allowbreak P_{10}\mathbin{\oplus}\allowbreak P_{18}\mathbin{\oplus}\allowbreak P_{19}\mathbin{\oplus}\allowbreak M_1$ \\
$M_3$ & $\ell_{\mathtt{75afc}}=U_1\mathbin{\oplus}\allowbreak U_2\mathbin{\oplus}\allowbreak U_3\mathbin{\oplus}\allowbreak U_4\mathbin{\oplus}\allowbreak U_5\mathbin{\oplus}\allowbreak U_6\mathbin{\oplus}\allowbreak P_{01}\mathbin{\oplus}\allowbreak P_{03}\mathbin{\oplus}\allowbreak P_{09}\mathbin{\oplus}\allowbreak P_{11}\mathbin{\oplus}\allowbreak P_{18}\mathbin{\oplus}\allowbreak P_{19}\mathbin{\oplus}\allowbreak M_1$ & $\ell_{\mathtt{2e080}}=U_6\mathbin{\oplus}\allowbreak P_{10}\mathbin{\oplus}\allowbreak P_{11}\mathbin{\oplus}\allowbreak P_{17}\mathbin{\oplus}\allowbreak P_{19}$ \\
$M_4$ & $\ell_{\mathtt{330ea}}=U_0\mathbin{\oplus}\allowbreak U_2\mathbin{\oplus}\allowbreak U_4\mathbin{\oplus}\allowbreak U_5\mathbin{\oplus}\allowbreak U_6\mathbin{\oplus}\allowbreak P_{09}\mathbin{\oplus}\allowbreak P_{10}\mathbin{\oplus}\allowbreak P_{18}\mathbin{\oplus}\allowbreak P_{19}$ & $\ell_{\mathtt{140000}}=M_1\mathbin{\oplus}\allowbreak M_3$ \\
$M_5$ & $\ell_{\mathtt{b5afc}}=U_1\mathbin{\oplus}\allowbreak U_2\mathbin{\oplus}\allowbreak U_3\mathbin{\oplus}\allowbreak U_4\mathbin{\oplus}\allowbreak U_5\mathbin{\oplus}\allowbreak U_6\mathbin{\oplus}\allowbreak P_{01}\mathbin{\oplus}\allowbreak P_{03}\mathbin{\oplus}\allowbreak P_{09}\mathbin{\oplus}\allowbreak P_{11}\mathbin{\oplus}\allowbreak P_{18}\mathbin{\oplus}\allowbreak P_{19}\mathbin{\oplus}\allowbreak M_2$ & $\ell_{\mathtt{2730ea}}=U_0\mathbin{\oplus}\allowbreak U_2\mathbin{\oplus}\allowbreak U_4\mathbin{\oplus}\allowbreak U_5\mathbin{\oplus}\allowbreak U_6\mathbin{\oplus}\allowbreak P_{09}\mathbin{\oplus}\allowbreak P_{10}\mathbin{\oplus}\allowbreak P_{18}\mathbin{\oplus}\allowbreak P_{19}\mathbin{\oplus}\allowbreak M_1\mathbin{\oplus}\allowbreak M_4$ \\
\end{longtable}
\endgroup

\subsection{Post-middle coordinates}
\begin{longtable}{@{}rc>{\raggedright\arraybackslash}p{0.72\textwidth}@{}}
\caption{The four post-middle coordinates in the same 23-element affine basis.}\label{tab:post-coordinates}\\
\toprule
coordinate & mask & affine expression \\
\midrule\endfirsthead
\toprule
coordinate & mask & affine expression \\
\midrule\endhead
\bottomrule\endfoot
$Z_0$ & \texttt{21d06a} & $U_0\mathbin{\oplus}\allowbreak U_2\mathbin{\oplus}\allowbreak U_4\mathbin{\oplus}\allowbreak U_5\mathbin{\oplus}\allowbreak P_{09}\mathbin{\oplus}\allowbreak P_{11}\mathbin{\oplus}\allowbreak P_{17}\mathbin{\oplus}\allowbreak P_{18}\mathbin{\oplus}\allowbreak M_4$ \\
$Z_1$ & \texttt{32e080} & $U_6\mathbin{\oplus}\allowbreak P_{10}\mathbin{\oplus}\allowbreak P_{11}\mathbin{\oplus}\allowbreak P_{17}\mathbin{\oplus}\allowbreak P_{19}\mathbin{\oplus}\allowbreak M_3\mathbin{\oplus}\allowbreak M_4$ \\
$Z_2$ & \texttt{49ecf6} & $U_0\mathbin{\oplus}\allowbreak U_1\mathbin{\oplus}\allowbreak U_3\mathbin{\oplus}\allowbreak U_4\mathbin{\oplus}\allowbreak U_5\mathbin{\oplus}\allowbreak U_6\mathbin{\oplus}\allowbreak P_{02}\mathbin{\oplus}\allowbreak P_{03}\mathbin{\oplus}\allowbreak P_{10}\mathbin{\oplus}\allowbreak P_{11}\mathbin{\oplus}\allowbreak P_{17}\mathbin{\oplus}\allowbreak P_{18}\mathbin{\oplus}\allowbreak M_2\mathbin{\oplus}\allowbreak M_5$ \\
$Z_3$ & \texttt{42b60a} & $U_0\mathbin{\oplus}\allowbreak U_2\mathbin{\oplus}\allowbreak P_{01}\mathbin{\oplus}\allowbreak P_{02}\mathbin{\oplus}\allowbreak P_{09}\mathbin{\oplus}\allowbreak P_{10}\mathbin{\oplus}\allowbreak P_{17}\mathbin{\oplus}\allowbreak P_{19}\mathbin{\oplus}\allowbreak M_5$ \\
\end{longtable}

\subsection{Output reconstruction}
For the output rows, append the late products in the order
\[
(P_{04},P_{05},P_{06},P_{07},P_{08},P_{12},P_{13},P_{14},P_{15},P_{16},P_{20},P_{21},P_{22},P_{23},P_{24})
\]
to the 23-element middle basis.
\begingroup\small
\begin{longtable}{@{}rc>{\raggedright\arraybackslash}p{0.72\textwidth}@{}}
\caption{Exact affine reconstruction of the eight output bits.}\label{tab:output-reconstruction}\\
\toprule
output & mask & affine expression \\
\midrule\endfirsthead
\toprule
output & mask & affine expression \\
\midrule\endhead
\bottomrule\endfoot
$S_0$ & \texttt{2491ba35ce} & $U_0\mathbin{\oplus}\allowbreak U_1\mathbin{\oplus}\allowbreak U_2\mathbin{\oplus}\allowbreak U_5\mathbin{\oplus}\allowbreak U_6\mathbin{\oplus}\allowbreak U_7\mathbin{\oplus}\allowbreak P_{02}\mathbin{\oplus}\allowbreak P_{09}\mathbin{\oplus}\allowbreak P_{10}\mathbin{\oplus}\allowbreak P_{19}\mathbin{\oplus}\allowbreak M_2\mathbin{\oplus}\allowbreak M_3\mathbin{\oplus}\allowbreak M_4\mathbin{\oplus}\allowbreak P_{04}\mathbin{\oplus}\allowbreak P_{05}\mathbin{\oplus}\allowbreak P_{12}\mathbin{\oplus}\allowbreak P_{15}\mathbin{\oplus}\allowbreak P_{21}\mathbin{\oplus}\allowbreak P_{24}$ \\
$S_1$ & \texttt{257bba3b95} & $1\mathbin{\oplus}\allowbreak U_1\mathbin{\oplus}\allowbreak U_3\mathbin{\oplus}\allowbreak U_6\mathbin{\oplus}\allowbreak U_7\mathbin{\oplus}\allowbreak P_{01}\mathbin{\oplus}\allowbreak P_{03}\mathbin{\oplus}\allowbreak P_{09}\mathbin{\oplus}\allowbreak P_{10}\mathbin{\oplus}\allowbreak P_{19}\mathbin{\oplus}\allowbreak M_2\mathbin{\oplus}\allowbreak M_3\mathbin{\oplus}\allowbreak M_4\mathbin{\oplus}\allowbreak P_{04}\mathbin{\oplus}\allowbreak P_{05}\mathbin{\oplus}\allowbreak P_{06}\mathbin{\oplus}\allowbreak P_{08}\mathbin{\oplus}\allowbreak P_{12}\mathbin{\oplus}\allowbreak P_{13}\mathbin{\oplus}\allowbreak P_{14}\mathbin{\oplus}\allowbreak P_{16}\mathbin{\oplus}\allowbreak P_{21}\mathbin{\oplus}\allowbreak P_{24}$ \\
$S_2$ & \texttt{2e6532ee01} & $1\mathbin{\oplus}\allowbreak P_{01}\mathbin{\oplus}\allowbreak P_{02}\mathbin{\oplus}\allowbreak P_{03}\mathbin{\oplus}\allowbreak P_{10}\mathbin{\oplus}\allowbreak P_{11}\mathbin{\oplus}\allowbreak P_{17}\mathbin{\oplus}\allowbreak P_{19}\mathbin{\oplus}\allowbreak M_3\mathbin{\oplus}\allowbreak M_4\mathbin{\oplus}\allowbreak P_{05}\mathbin{\oplus}\allowbreak P_{07}\mathbin{\oplus}\allowbreak P_{13}\mathbin{\oplus}\allowbreak P_{14}\mathbin{\oplus}\allowbreak P_{20}\mathbin{\oplus}\allowbreak P_{21}\mathbin{\oplus}\allowbreak P_{22}\mathbin{\oplus}\allowbreak P_{24}$ \\
$S_3$ & \texttt{19dbba4b98} & $U_2\mathbin{\oplus}\allowbreak U_3\mathbin{\oplus}\allowbreak U_6\mathbin{\oplus}\allowbreak U_7\mathbin{\oplus}\allowbreak P_{01}\mathbin{\oplus}\allowbreak P_{03}\mathbin{\oplus}\allowbreak P_{11}\mathbin{\oplus}\allowbreak P_{19}\mathbin{\oplus}\allowbreak M_2\mathbin{\oplus}\allowbreak M_3\mathbin{\oplus}\allowbreak M_4\mathbin{\oplus}\allowbreak P_{04}\mathbin{\oplus}\allowbreak P_{05}\mathbin{\oplus}\allowbreak P_{06}\mathbin{\oplus}\allowbreak P_{08}\mathbin{\oplus}\allowbreak P_{12}\mathbin{\oplus}\allowbreak P_{14}\mathbin{\oplus}\allowbreak P_{15}\mathbin{\oplus}\allowbreak P_{16}\mathbin{\oplus}\allowbreak P_{22}\mathbin{\oplus}\allowbreak P_{23}$ \\
$S_4$ & \texttt{189eba357c} & $U_1\mathbin{\oplus}\allowbreak U_2\mathbin{\oplus}\allowbreak U_3\mathbin{\oplus}\allowbreak U_4\mathbin{\oplus}\allowbreak U_5\mathbin{\oplus}\allowbreak U_7\mathbin{\oplus}\allowbreak P_{02}\mathbin{\oplus}\allowbreak P_{09}\mathbin{\oplus}\allowbreak P_{10}\mathbin{\oplus}\allowbreak P_{19}\mathbin{\oplus}\allowbreak M_2\mathbin{\oplus}\allowbreak M_3\mathbin{\oplus}\allowbreak M_4\mathbin{\oplus}\allowbreak P_{04}\mathbin{\oplus}\allowbreak P_{06}\mathbin{\oplus}\allowbreak P_{07}\mathbin{\oplus}\allowbreak P_{08}\mathbin{\oplus}\allowbreak P_{12}\mathbin{\oplus}\allowbreak P_{15}\mathbin{\oplus}\allowbreak P_{22}\mathbin{\oplus}\allowbreak P_{23}$ \\
$S_5$ & \texttt{0c0332e152} & $U_0\mathbin{\oplus}\allowbreak U_3\mathbin{\oplus}\allowbreak U_5\mathbin{\oplus}\allowbreak U_7\mathbin{\oplus}\allowbreak P_{10}\mathbin{\oplus}\allowbreak P_{11}\mathbin{\oplus}\allowbreak P_{17}\mathbin{\oplus}\allowbreak P_{19}\mathbin{\oplus}\allowbreak M_3\mathbin{\oplus}\allowbreak M_4\mathbin{\oplus}\allowbreak P_{05}\mathbin{\oplus}\allowbreak P_{06}\mathbin{\oplus}\allowbreak P_{21}\mathbin{\oplus}\allowbreak P_{22}$ \\
$S_6$ & \texttt{0db00b2bf9} & $1\mathbin{\oplus}\allowbreak U_2\mathbin{\oplus}\allowbreak U_3\mathbin{\oplus}\allowbreak U_4\mathbin{\oplus}\allowbreak U_5\mathbin{\oplus}\allowbreak U_6\mathbin{\oplus}\allowbreak U_7\mathbin{\oplus}\allowbreak P_{01}\mathbin{\oplus}\allowbreak P_{03}\mathbin{\oplus}\allowbreak P_{10}\mathbin{\oplus}\allowbreak P_{18}\mathbin{\oplus}\allowbreak P_{19}\mathbin{\oplus}\allowbreak M_2\mathbin{\oplus}\allowbreak P_{12}\mathbin{\oplus}\allowbreak P_{13}\mathbin{\oplus}\allowbreak P_{15}\mathbin{\oplus}\allowbreak P_{16}\mathbin{\oplus}\allowbreak P_{21}\mathbin{\oplus}\allowbreak P_{22}$ \\
$S_7$ & \texttt{19b508a44b} & $1\mathbin{\oplus}\allowbreak U_0\mathbin{\oplus}\allowbreak U_2\mathbin{\oplus}\allowbreak U_5\mathbin{\oplus}\allowbreak P_{02}\mathbin{\oplus}\allowbreak P_{10}\mathbin{\oplus}\allowbreak P_{17}\mathbin{\oplus}\allowbreak M_2\mathbin{\oplus}\allowbreak P_{05}\mathbin{\oplus}\allowbreak P_{07}\mathbin{\oplus}\allowbreak P_{12}\mathbin{\oplus}\allowbreak P_{13}\mathbin{\oplus}\allowbreak P_{15}\mathbin{\oplus}\allowbreak P_{16}\mathbin{\oplus}\allowbreak P_{22}\mathbin{\oplus}\allowbreak P_{23}$ \\
\end{longtable}
\endgroup


\subsection{Formal target and reconstruction masks}
\begin{longtable}{@{}r>{\raggedright\arraybackslash}p{0.27\textwidth}ll@{}}
\caption{Formal quadratic reconstruction certificate derived from the NIST SLP. The first four rows form the pre-middle subspace.}\label{tab:targets}\\
\toprule
row & first independent requirement & target signature & selected products \\
\midrule
\endfirsthead
\toprule
row & first independent requirement & target signature & selected products \\
\midrule
\endhead
\bottomrule
\endfoot
01 ($Q_{\mathrm{pre}}$) & first factor of $M_1$ & \texttt{000000003c281a058} & \texttt{030606} \\
02 ($Q_{\mathrm{pre}}$) & second factor of $M_1$ & \texttt{000000001c381e050} & \texttt{030500} \\
03 ($Q_{\mathrm{pre}}$) & first factor of $M_2$ & \texttt{00000000281416030} & \texttt{050303} \\
04 ($Q_{\mathrm{pre}}$) & second factor of $M_2$ & \texttt{00000000241c0a078} & \texttt{060300} \\
05 ($Q$) & output $S_0$ & \texttt{007e7943c281a0580} & \texttt{97481f} \\
06 ($Q$) & output $S_1$ & \texttt{00fa6963c281a0580} & \texttt{97b8b8} \\
07 ($Q$) & output $S_2$ & \texttt{005f31a241c0a0780} & \texttt{b83057} \\
08 ($Q$) & output $S_3$ & \texttt{00841023c281a0580} & \texttt{67efb8} \\
09 ($Q$) & output $S_4$ & \texttt{004c0863c281a0580} & \texttt{6748ef} \\
10 ($Q$) & output $S_5$ & \texttt{002148e0c081c0480} & \texttt{300030} \\
11 ($Q$) & output $S_6$ & \texttt{00e950a1c381e0500} & \texttt{30df00} \\
12 ($Q$) & output $S_7$ & \texttt{00c81841c381e0500} & \texttt{67df57} \\
\end{longtable}


\section{Complete finite evaluation and normalized graph}
\label{app:finite}

\begingroup\scriptsize
\setlength{\tabcolsep}{2.7pt}
\begin{longtable}{c*{16}{c}}
\caption{Complete evaluation of the staged 29-AND certificate. Row and column headers are hexadecimal input nibbles; each body entry is the hexadecimal output byte. This table is entrywise identical to FIPS 197, Table 4.}\label{tab:evaluation}\\
\toprule
 & \texttt{0} & \texttt{1} & \texttt{2} & \texttt{3} & \texttt{4} & \texttt{5} & \texttt{6} & \texttt{7} & \texttt{8} & \texttt{9} & \texttt{a} & \texttt{b} & \texttt{c} & \texttt{d} & \texttt{e} & \texttt{f} \\
\midrule
\endfirsthead
\toprule
 & \texttt{0} & \texttt{1} & \texttt{2} & \texttt{3} & \texttt{4} & \texttt{5} & \texttt{6} & \texttt{7} & \texttt{8} & \texttt{9} & \texttt{a} & \texttt{b} & \texttt{c} & \texttt{d} & \texttt{e} & \texttt{f} \\
\midrule
\endhead
\bottomrule
\endfoot
\texttt{0} & \texttt{63} & \texttt{7c} & \texttt{77} & \texttt{7b} & \texttt{f2} & \texttt{6b} & \texttt{6f} & \texttt{c5} & \texttt{30} & \texttt{01} & \texttt{67} & \texttt{2b} & \texttt{fe} & \texttt{d7} & \texttt{ab} & \texttt{76} \\
\texttt{1} & \texttt{ca} & \texttt{82} & \texttt{c9} & \texttt{7d} & \texttt{fa} & \texttt{59} & \texttt{47} & \texttt{f0} & \texttt{ad} & \texttt{d4} & \texttt{a2} & \texttt{af} & \texttt{9c} & \texttt{a4} & \texttt{72} & \texttt{c0} \\
\texttt{2} & \texttt{b7} & \texttt{fd} & \texttt{93} & \texttt{26} & \texttt{36} & \texttt{3f} & \texttt{f7} & \texttt{cc} & \texttt{34} & \texttt{a5} & \texttt{e5} & \texttt{f1} & \texttt{71} & \texttt{d8} & \texttt{31} & \texttt{15} \\
\texttt{3} & \texttt{04} & \texttt{c7} & \texttt{23} & \texttt{c3} & \texttt{18} & \texttt{96} & \texttt{05} & \texttt{9a} & \texttt{07} & \texttt{12} & \texttt{80} & \texttt{e2} & \texttt{eb} & \texttt{27} & \texttt{b2} & \texttt{75} \\
\texttt{4} & \texttt{09} & \texttt{83} & \texttt{2c} & \texttt{1a} & \texttt{1b} & \texttt{6e} & \texttt{5a} & \texttt{a0} & \texttt{52} & \texttt{3b} & \texttt{d6} & \texttt{b3} & \texttt{29} & \texttt{e3} & \texttt{2f} & \texttt{84} \\
\texttt{5} & \texttt{53} & \texttt{d1} & \texttt{00} & \texttt{ed} & \texttt{20} & \texttt{fc} & \texttt{b1} & \texttt{5b} & \texttt{6a} & \texttt{cb} & \texttt{be} & \texttt{39} & \texttt{4a} & \texttt{4c} & \texttt{58} & \texttt{cf} \\
\texttt{6} & \texttt{d0} & \texttt{ef} & \texttt{aa} & \texttt{fb} & \texttt{43} & \texttt{4d} & \texttt{33} & \texttt{85} & \texttt{45} & \texttt{f9} & \texttt{02} & \texttt{7f} & \texttt{50} & \texttt{3c} & \texttt{9f} & \texttt{a8} \\
\texttt{7} & \texttt{51} & \texttt{a3} & \texttt{40} & \texttt{8f} & \texttt{92} & \texttt{9d} & \texttt{38} & \texttt{f5} & \texttt{bc} & \texttt{b6} & \texttt{da} & \texttt{21} & \texttt{10} & \texttt{ff} & \texttt{f3} & \texttt{d2} \\
\texttt{8} & \texttt{cd} & \texttt{0c} & \texttt{13} & \texttt{ec} & \texttt{5f} & \texttt{97} & \texttt{44} & \texttt{17} & \texttt{c4} & \texttt{a7} & \texttt{7e} & \texttt{3d} & \texttt{64} & \texttt{5d} & \texttt{19} & \texttt{73} \\
\texttt{9} & \texttt{60} & \texttt{81} & \texttt{4f} & \texttt{dc} & \texttt{22} & \texttt{2a} & \texttt{90} & \texttt{88} & \texttt{46} & \texttt{ee} & \texttt{b8} & \texttt{14} & \texttt{de} & \texttt{5e} & \texttt{0b} & \texttt{db} \\
\texttt{a} & \texttt{e0} & \texttt{32} & \texttt{3a} & \texttt{0a} & \texttt{49} & \texttt{06} & \texttt{24} & \texttt{5c} & \texttt{c2} & \texttt{d3} & \texttt{ac} & \texttt{62} & \texttt{91} & \texttt{95} & \texttt{e4} & \texttt{79} \\
\texttt{b} & \texttt{e7} & \texttt{c8} & \texttt{37} & \texttt{6d} & \texttt{8d} & \texttt{d5} & \texttt{4e} & \texttt{a9} & \texttt{6c} & \texttt{56} & \texttt{f4} & \texttt{ea} & \texttt{65} & \texttt{7a} & \texttt{ae} & \texttt{08} \\
\texttt{c} & \texttt{ba} & \texttt{78} & \texttt{25} & \texttt{2e} & \texttt{1c} & \texttt{a6} & \texttt{b4} & \texttt{c6} & \texttt{e8} & \texttt{dd} & \texttt{74} & \texttt{1f} & \texttt{4b} & \texttt{bd} & \texttt{8b} & \texttt{8a} \\
\texttt{d} & \texttt{70} & \texttt{3e} & \texttt{b5} & \texttt{66} & \texttt{48} & \texttt{03} & \texttt{f6} & \texttt{0e} & \texttt{61} & \texttt{35} & \texttt{57} & \texttt{b9} & \texttt{86} & \texttt{c1} & \texttt{1d} & \texttt{9e} \\
\texttt{e} & \texttt{e1} & \texttt{f8} & \texttt{98} & \texttt{11} & \texttt{69} & \texttt{d9} & \texttt{8e} & \texttt{94} & \texttt{9b} & \texttt{1e} & \texttt{87} & \texttt{e9} & \texttt{ce} & \texttt{55} & \texttt{28} & \texttt{df} \\
\texttt{f} & \texttt{8c} & \texttt{a1} & \texttt{89} & \texttt{0d} & \texttt{bf} & \texttt{e6} & \texttt{42} & \texttt{68} & \texttt{41} & \texttt{99} & \texttt{2d} & \texttt{0f} & \texttt{b0} & \texttt{54} & \texttt{bb} & \texttt{16} \\
\end{longtable}
\endgroup


\subsection{Normalized XOR--AND graph}

The following certificate absorbs all affine gates into masks.  Mask bit 0 is the constant 1, mask
bits 1--8 are \(U_0,\ldots,U_7\), and each later basis bit is the output of one preceding AND row.
It therefore gives a compact, implementation-independent description of the 29 nonlinear gates and
the eight final affine output masks.

\begin{lstlisting}[language=bash,caption={Normalized 29-AND XAG certificate.},label={lst:xag}]
AES29-XAG 1
AND_COUNT 29
AND A01 1e4 10e
AND A02 10c 18c
AND A03 e8 82
AND A04 1b2 1ce
AND A05 b2 d2
AND A06 100 11c
AND A07 56 c0
AND A08 1be 15e
AND A09 1e8 19e
AND A10 1ecf6 1d06a
AND A11 2b60a 730ea
AND A12 75afc 2e080
AND A13 330ea 140000
AND A14 b5afc 2730ea
AND A15 42b7ec b5a94
AND A16 42b76e b5b9a
AND A17 49ed92 b5afe
AND A18 42b662 b5b72
AND A19 49ec9e b5b1a
AND A20 21d0a4 32e1e0
AND A21 21d1b8 32e02e
AND A22 133138 32e1fc
AND A23 21d10a 32e12e
AND A24 13318a 32e04e
AND A25 636748 39bb74
AND A26 6366d6 39bbb4
AND A27 5adcaa 39bb02
AND A28 636768 39ba5c
AND A29 5add14 39bb54
OUTPUT_COUNT 8
OUTPUT S0 2491ba35ce
OUTPUT S1 257bba3b95
OUTPUT S2 2e6532ee01
OUTPUT S3 19dbba4b98
OUTPUT S4 189eba357c
OUTPUT S5 c0332e152
OUTPUT S6 db00b2bf9
OUTPUT S7 19b508a44b
END
\end{lstlisting}


\section{Repository artifacts and standalone verification}
\label{app:artifacts}

This final appendix identifies the executable artifacts corresponding to the mathematical objects
above.  Internal repository paths are confined to this appendix.

\begin{longtable}{@{}>{\raggedright\arraybackslash}p{0.43\textwidth}>{\raggedright\arraybackslash}p{0.49\textwidth}@{}}
\caption{Principal released artifacts and fixed SHA-256 digests.}\label{tab:artifacts}\\
\toprule
artifact & digest or role \\
\midrule
\endfirsthead
\toprule
artifact & digest or role \\
\midrule
\endhead
\bottomrule
\endfoot
\path{circuits/aes-sbox-fwd-g228-a29-d35-ad6.slp} & \primaryhash \\
\path{circuits/aes-sbox-fwd-g455-a29-d35-ad6-transparent.slp} & \transparenthash \\
\path{certificates/aes29.xag} & \xaghash \\
\path{certificates/tower24_witness.txt} & \towerhash \\
\path{certificates/staged29_witness.txt} & \stagedhash \\
\path{baseline/aes-sbox-fwd-g113-a32-d27-ad6.slp} & exact bytes of the official NIST SLP; \baselinehash \\
\path{verify/verify_staged29.py} & strict parser and exhaustive evaluator for the complete staged certificate \\
\path{verify/minimal_verify.py} & strict standalone SLP parser, FIPS algebraic/table cross-check, exhaustive evaluator, tally, and depth checker \\
\path{verify/verify_tower24.py} & independent NIST target derivation, formal 4096-point bridge, and staged 256-point check \\
\path{verify/run_all.sh} & redundant full verification suite \\
\end{longtable}

The shortest checks of the two direct certificates are
\begin{lstlisting}[language=bash]
python3 verify/verify_staged29.py
python3 verify/minimal_verify.py
\end{lstlisting}
The complete redundant audit is
\begin{lstlisting}[language=bash]
./verify/run_all.sh
\end{lstlisting}
The canonical paper source is \path{article/article.tex}.  The release process first generates the
included tables, compiles that TeX source to \path{article/article.pdf}, and only after the PDF
exists renders the same expanded TeX source as the mathematical GitHub Pages document
\path{docs/index.html}.  The HTML is therefore a second rendering of the paper, not a separate
informal article.  Public copies of the circuit, certificates, verifier, and PDF are placed in
\path{docs/files/} and checked byte for byte.

\subsection{Standalone exhaustive implementation verifier}

The following listing uses only the Python standard library.  It contains a literal copy of the FIPS
S-box table, independently computes the algebraic definition, strictly parses the complete released
SLP container, checks exact counts and depths, and evaluates all 256 inputs.

\begin{lstlisting}[language=Python,caption={Standalone exhaustive verifier.},label={lst:minimal}]
#!/usr/bin/env python3
"""Standalone exhaustive verifier for the released 29-AND AES S-box SLP.

Only the Python standard library is used.  The checker strictly parses the
complete circuit container, verifies all data dependencies and exact gate
counts, independently checks the literal FIPS 197 table against the algebraic
S-box definition, and evaluates the circuit on every one of the 256 inputs.
U0 and S0 are the most significant bits.
"""
from __future__ import annotations

from collections import Counter
from pathlib import Path
import hashlib
import re
import sys

ROOT = Path(__file__).resolve().parents[1]
DEFAULT = ROOT / "circuits" / "aes-sbox-fwd-g228-a29-d35-ad6.slp"
DEFAULT_SHA256 = "f41861c4b15bc78840708a5da2fab6ca9dd204c1e3e572cae49912c713bf18d3"
FIPS_SBOX = (
    0x63, 0x7c, 0x77, 0x7b, 0xf2, 0x6b, 0x6f, 0xc5, 0x30, 0x01, 0x67, 0x2b, 0xfe, 0xd7, 0xab, 0x76,
    0xca, 0x82, 0xc9, 0x7d, 0xfa, 0x59, 0x47, 0xf0, 0xad, 0xd4, 0xa2, 0xaf, 0x9c, 0xa4, 0x72, 0xc0,
    0xb7, 0xfd, 0x93, 0x26, 0x36, 0x3f, 0xf7, 0xcc, 0x34, 0xa5, 0xe5, 0xf1, 0x71, 0xd8, 0x31, 0x15,
    0x04, 0xc7, 0x23, 0xc3, 0x18, 0x96, 0x05, 0x9a, 0x07, 0x12, 0x80, 0xe2, 0xeb, 0x27, 0xb2, 0x75,
    0x09, 0x83, 0x2c, 0x1a, 0x1b, 0x6e, 0x5a, 0xa0, 0x52, 0x3b, 0xd6, 0xb3, 0x29, 0xe3, 0x2f, 0x84,
    0x53, 0xd1, 0x00, 0xed, 0x20, 0xfc, 0xb1, 0x5b, 0x6a, 0xcb, 0xbe, 0x39, 0x4a, 0x4c, 0x58, 0xcf,
    0xd0, 0xef, 0xaa, 0xfb, 0x43, 0x4d, 0x33, 0x85, 0x45, 0xf9, 0x02, 0x7f, 0x50, 0x3c, 0x9f, 0xa8,
    0x51, 0xa3, 0x40, 0x8f, 0x92, 0x9d, 0x38, 0xf5, 0xbc, 0xb6, 0xda, 0x21, 0x10, 0xff, 0xf3, 0xd2,
    0xcd, 0x0c, 0x13, 0xec, 0x5f, 0x97, 0x44, 0x17, 0xc4, 0xa7, 0x7e, 0x3d, 0x64, 0x5d, 0x19, 0x73,
    0x60, 0x81, 0x4f, 0xdc, 0x22, 0x2a, 0x90, 0x88, 0x46, 0xee, 0xb8, 0x14, 0xde, 0x5e, 0x0b, 0xdb,
    0xe0, 0x32, 0x3a, 0x0a, 0x49, 0x06, 0x24, 0x5c, 0xc2, 0xd3, 0xac, 0x62, 0x91, 0x95, 0xe4, 0x79,
    0xe7, 0xc8, 0x37, 0x6d, 0x8d, 0xd5, 0x4e, 0xa9, 0x6c, 0x56, 0xf4, 0xea, 0x65, 0x7a, 0xae, 0x08,
    0xba, 0x78, 0x25, 0x2e, 0x1c, 0xa6, 0xb4, 0xc6, 0xe8, 0xdd, 0x74, 0x1f, 0x4b, 0xbd, 0x8b, 0x8a,
    0x70, 0x3e, 0xb5, 0x66, 0x48, 0x03, 0xf6, 0x0e, 0x61, 0x35, 0x57, 0xb9, 0x86, 0xc1, 0x1d, 0x9e,
    0xe1, 0xf8, 0x98, 0x11, 0x69, 0xd9, 0x8e, 0x94, 0x9b, 0x1e, 0x87, 0xe9, 0xce, 0x55, 0x28, 0xdf,
    0x8c, 0xa1, 0x89, 0x0d, 0xbf, 0xe6, 0x42, 0x68, 0x41, 0x99, 0x2d, 0x0f, 0xb0, 0x54, 0xbb, 0x16,
)


def require(condition: bool, message: str) -> None:
    if not condition:
        raise SystemExit(f"FAIL: {message}")


def gf256_mul(a: int, b: int) -> int:
    result = 0
    for _ in range(8):
        if b & 1:
            result ^= a
        a = ((a << 1) & 0xFF) ^ (0x1B if a & 0x80 else 0)
        b >>= 1
    return result


def aes_sbox_algebraic(x: int) -> int:
    inverse = 0
    if x:
        inverse, base, exponent = 1, x, 254
        while exponent:
            if exponent & 1:
                inverse = gf256_mul(inverse, base)
            base = gf256_mul(base, base)
            exponent >>= 1
    output = 0
    for i in range(8):
        bit = (
            ((inverse >> i) & 1)
            ^ ((inverse >> ((i + 4) & 7)) & 1)
            ^ ((inverse >> ((i + 5) & 7)) & 1)
            ^ ((inverse >> ((i + 6) & 7)) & 1)
            ^ ((inverse >> ((i + 7) & 7)) & 1)
            ^ ((0x63 >> i) & 1)
        )
        output |= bit << i
    return output


def parse(path: Path) -> tuple[list[tuple[str, str, str, str | None]], tuple[str, ...], tuple[str, ...]]:
    operations: list[tuple[str, str, str, str | None]] = []
    inputs = tuple(f"U{i}" for i in range(8))
    outputs = tuple(f"S{i}" for i in range(8))
    state = "start"
    saw_inputs = saw_outputs = saw_internal = saw_syntax = False
    internal_count: int | None = None

    for lineno, raw in enumerate(path.read_text(encoding="utf-8").splitlines(), 1):
        line = raw.split("#", 1)[0].strip()
        if not line:
            continue

        if state == "start":
            require(re.fullmatch(r"begin circuit [A-Za-z0-9_-]+", line) is not None,
                    f"expected begin circuit at line {lineno}")
            state = "header"
            continue

        if state == "header":
            if line == "Inputs: U0:U7":
                require(not saw_inputs, f"duplicate input declaration at line {lineno}")
                saw_inputs = True
                continue
            if line == "Outputs: S0:S7":
                require(saw_inputs and not saw_outputs, f"misordered or duplicate output declaration at line {lineno}")
                saw_outputs = True
                continue
            match = re.fullmatch(r"Internal: t1:t([1-9][0-9]*)", line)
            if match:
                require(saw_outputs and not saw_internal, f"misordered or duplicate internal declaration at line {lineno}")
                internal_count = int(match.group(1))
                saw_internal = True
                continue
            if line == "GateSyntax: GateName Output Inputs":
                require(saw_internal and not saw_syntax, f"misordered or duplicate gate syntax at line {lineno}")
                saw_syntax = True
                continue
            if line == "begin SLP":
                require(saw_inputs and saw_outputs and saw_internal and saw_syntax,
                        f"incomplete circuit header before line {lineno}")
                state = "slp"
                continue
            require(False, f"unexpected circuit-header line {lineno}: {line}")

        if state == "slp":
            if line == "end SLP":
                state = "after_slp"
                continue
            fields = line.split()
            if len(fields) == 4 and fields[0] in {"XOR", "AND"}:
                operations.append((fields[0], fields[1], fields[2], fields[3]))
            elif len(fields) == 3 and fields[0] == "NOT":
                operations.append((fields[0], fields[1], fields[2], None))
            else:
                require(False, f"illegal instruction at line {lineno}: {line}")
            continue

        if state == "after_slp":
            require(line == "end circuit", f"expected end circuit at line {lineno}")
            state = "done"
            continue

        require(False, f"trailing non-comment data at line {lineno}: {line}")

    require(state == "done", "missing or unbalanced circuit/SLP delimiters")
    require(internal_count is not None, "missing internal-wire declaration")
    require(bool(operations), "empty SLP")

    defined = set(inputs)
    temporary_outputs: set[str] = set()
    for kind, output, a, b in operations:
        require(output not in defined, f"wire redefinition: {output}")
        require(a in defined, f"use before definition: {a}")
        require(b is None or b in defined, f"use before definition: {b}")
        require(re.fullmatch(r"t[1-9][0-9]*|S[0-7]", output) is not None,
                f"undeclared output-wire name: {output}")
        if output.startswith("t"):
            number = int(output[1:])
            require(1 <= number <= internal_count, f"temporary outside declared range: {output}")
            temporary_outputs.add(output)
        defined.add(output)

    declared_temporaries = {f"t{i}" for i in range(1, internal_count + 1)}
    require(temporary_outputs == declared_temporaries, "declared and defined temporary-wire sets differ")
    require(all(name in defined for name in outputs), "missing output wire")
    return operations, inputs, outputs


def evaluate(operations: list[tuple[str, str, str, str | None]], inputs: tuple[str, ...], outputs: tuple[str, ...], x: int) -> int:
    wire = {name: (x >> (7 - i)) & 1 for i, name in enumerate(inputs)}
    for kind, output, a, b in operations:
        if kind == "XOR":
            wire[output] = wire[a] ^ wire[b]  # type: ignore[index]
        elif kind == "AND":
            wire[output] = wire[a] & wire[b]  # type: ignore[index]
        else:
            wire[output] = wire[a] ^ 1
    return sum(wire[name] << (7 - i) for i, name in enumerate(outputs))


def depths(operations: list[tuple[str, str, str, str | None]], inputs: tuple[str, ...], outputs: tuple[str, ...]) -> tuple[int, int]:
    gate_depth = {name: 0 for name in inputs}
    and_depth = dict(gate_depth)
    for kind, output, a, b in operations:
        operands = (a,) if b is None else (a, b)
        gate_depth[output] = 1 + max(gate_depth[name] for name in operands)
        and_depth[output] = max(and_depth[name] for name in operands) + (1 if kind == "AND" else 0)
    return max(gate_depth[name] for name in outputs), max(and_depth[name] for name in outputs)


def main() -> None:
    require(len(sys.argv) <= 2, "usage: minimal_verify.py [circuit.slp]")
    path = Path(sys.argv[1]) if len(sys.argv) == 2 else DEFAULT
    operations, inputs, outputs = parse(path)

    counts = Counter(kind for kind, _, _, _ in operations)
    expected_counts = Counter({"AND": 29, "XOR": 195, "NOT": 4})
    require(counts == expected_counts, f"unexpected operation tally: {dict(counts)}")
    require(len(operations) == 228, f"expected exactly 228 instructions, found {len(operations)}")
    require(depths(operations, inputs, outputs) == (35, 6), "unexpected ordinary depth or AND-depth")

    for x, table_value in enumerate(FIPS_SBOX):
        algebraic = aes_sbox_algebraic(x)
        require(algebraic == table_value,
                f"FIPS algebraic definition and literal table disagree at {x:02x}: {algebraic:02x} != {table_value:02x}")
        obtained = evaluate(operations, inputs, outputs, x)
        require(obtained == table_value,
                f"input {x:02x}: obtained {obtained:02x}, expected {table_value:02x}")

    digest = hashlib.sha256(path.read_bytes()).hexdigest()
    if path.resolve() == DEFAULT.resolve():
        require(digest == DEFAULT_SHA256, f"default certificate SHA-256 mismatch: {digest}")
    print(
        "PASS: "
        f"{len(operations)} instructions; {counts['AND']} AND, {counts['XOR']} XOR, "
        f"{counts['NOT']} NOT; depth 35; AND-depth 6; all 256 FIPS entries; sha256 {digest}"
    )


if __name__ == "__main__":
    main()
\end{lstlisting}


\printbibliography

\end{document}
